\chapter{부도확률 $p$는 어떻게 만들어지고, 왜 믿을 수 있나}

이 장은 팀에서 가장 자주 나온 질문에 답한다. ``$p$가 모델에서 나온다는데, 정확히 어떻게 계산되나? 그 숫자를 왜 믿어야 하나?'' 답을 먼저 쓰면 이렇다.

\begin{keybox}[이 장의 답]
\begin{itemize}[leftmargin=*]
\item $p$는 ``이 대출이 만기까지 가는 동안 결국 부도로 끝날 확률''이다. LightGBM이라는 모델이 과거 대출의 결과를 보고 배운 규칙으로 계산한다.
\item 모델은 수백 개의 작은 결정나무를 차례로 쌓는다. 각 나무는 ``앞 나무들이 틀린 만큼''을 고치고, 나무들의 점수를 모두 더한 뒤 S자 함수로 0--1 사이 확률로 바꾼 것이 $p$다.
\item $p$의 \textbf{순서}는 믿을 수 있다. $p$를 10등분하면 실제 부도율이 3.8\%에서 30.6\%까지 차례로 오르고, 같은 등급 안에서도 그렇다. 우연이 아니라는 검증(정답 섞기, 시드, 다른 모델)도 통과했다.
\item $p$의 \textbf{크기}는 그대로 믿으면 안 된다. 실제보다 낮게 나오는 경향이 있다. 그래서 우리는 $p$를 순위를 매기는 데 쓰고, 그 순위가 실제로 돈을 버는지는 IRR로 따로 채점한다.
\end{itemize}
\Intuit{$p$는 정확한 온도계가 아니라 믿을 만한 줄 세우기다. 누가 더 위험한지는 잘 맞히고, 몇 \%인지는 대략 맞힌다.}
\end{keybox}

\section{$p$는 무엇인가}
\paragraph{정답(라벨).} 이미 끝난 대출에 대해 부도(Charged Off, Default)면 1, 다 갚았으면(Fully Paid) 0이다. 진행 중인 대출은 결과를 모르므로 학습에도 채점에도 쓰지 않는다(3장).

\paragraph{입력(특성).} 대출을 승인하던 날 알 수 있던 정보 101개다. 금리, 등급, 대출액, 소득, 부채비율(DTI), FICO, 근속 연수, 거주 주, 최근 계좌 개설 수, 신용 이력 길이 등이다. 결과를 담은 컬럼(받은 총액, 마지막 납부일, 대출 후 갱신된 신용점수)은 쓰지 않는다(1장).

\paragraph{뜻.} $p = 0.15$는 ``이 대출과 비슷한 대출 100개 중 약 15개가 결국 부도났다''는 뜻이다. 1년 안의 부도확률이 아니라 대출 생애 전체 기준이다.

\section{가장 단순한 모델부터: 비율과 로지스틱}
가장 단순한 $p$는 ``학습 데이터의 부도 비율''이다. 2016년 모델의 학습 데이터에서는 약 17\%였다. 모든 대출에 17\%를 주면 아무것도 구분하지 못한다.

다음 단계는 로지스틱 회귀다. 특성의 가중합 $z=\beta_0+\beta_1x_1+\dots$를 만들고 $p=1/(1+e^{-z})$로 바꾼다. $z$는 ``로그 오즈'' $\ln\frac{p}{1-p}$이며, 계수 $\beta_j$는 ``$x_j$가 1 늘면 로그 오즈가 $\beta_j$ 는다''로 읽는다. 문제는 직선 하나로 모든 관계를 표현해야 한다는 점이다. ``소득이 낮으면서 부채비율도 높을 때만 위험이 급증한다'' 같은 조합은 잘 잡지 못한다.

\section{결정나무: 질문으로 나누기}
결정나무는 ``FICO(미국 개인 신용점수)가 690보다 낮은가?''(기준점은 설명용 예시이고, 실제 기준점은 모델이 데이터에서 찾는다) 같은 질문으로 대출을 둘로 나누고, 각 칸의 부도율을 예측으로 쓴다. 질문을 여러 번 하면 조합 효과를 자연스럽게 잡는다. 결측값도 ``값이 없는 쪽''으로 보내는 법을 스스로 배운다(4장). 하지만 나무 하나는 학습 데이터의 우연에 민감해서 불안정하다.

\section{그래디언트 부스팅: 틀린 만큼 고치며 쌓기}
LightGBM은 나무를 \textbf{차례로} 쌓는다. 각 나무는 처음부터 다시 배우지 않고, 앞 나무들이 아직 틀리고 있는 부분만 고친다. 무엇이 ``틀린 부분''인지는 잔차 $y-p$로 잰다. 부도 났는데($y=1$) $p$가 낮으면 잔차가 크게 양수이고, 그 대출의 점수를 올려야 한다는 뜻이다.

\begin{workout}[대출 6건으로 부스팅 두 번]
\Use{로그 오즈, 잔차, 뉴턴 단계(아래 설명), 학습률}
\Step{1} 예시 데이터(숫자는 예시): 대출 6건, 특성 두 개(FICO 낮음 여부, 최근 계좌 개설 많음 여부), 부도 2건.
\begin{center}\small
\begin{tabular}{@{}lcccccc@{}}\toprule
대출 & 1 & 2 & 3 & 4 & 5 & 6\\\midrule
FICO 낮음 & 1 & 1 & 1 & 0 & 0 & 0\\
계좌 개설 많음 & 1 & 0 & 1 & 1 & 0 & 0\\
부도 $y$ & 1 & 0 & 1 & 0 & 0 & 0\\\bottomrule
\end{tabular}\end{center}
\Step{2} 출발점: 전체 부도율 $2/6$의 로그 오즈 $\ln(1/2)=-0.693$을 모든 대출의 점수 $z$로 둔다. $p=0.333$. 로그 손실 0.637. \Concept{로그 오즈}
\Step{3} 나무 1(질문: FICO 낮음?). 잔차 $y-p$는 부도 대출 $+0.667$, 정상 대출 $-0.333$. ``FICO 낮음'' 칸(대출 1--3)의 잔차 합 $+1.0$, 곡률 합 $\sum p(1-p)=0.667$이므로 잎 값 $1.0/0.667=+1.5$. 반대 칸은 $-1.5$. \Concept{뉴턴 단계}
\Step{4} 학습률 0.5를 곱해 점수에 더한다: FICO 낮음 칸 $z=-0.693+0.75=0.057$, 나머지 $-1.443$. 확률은 0.514와 0.191, 로그 손실 0.448.
\Step{5} 나무 2(질문: 계좌 개설 많음?). 새 잔차로 같은 계산: ``많음'' 칸 잎 값 $+1.193$, ``적음'' 칸 $-1.604$. 학습률을 곱해 $+0.597$, $-0.802$를 더한다.
\Step{6} 최종 $p$: 대출 1·\allowbreak{}3(부도) 0.658, 대출 2 0.322, 대출 4 0.300, 대출 5·\allowbreak{}6 0.096. 로그 손실 0.297.
\Ans{나무 두 개만으로 부도 대출의 $p$가 0.33에서 0.66으로 오르고, 안전한 대출은 0.10으로 내려갔다. 실제 모델은 이런 나무를 수백 개 쌓는다.}
\Verify{\ft{docs/notes/lecture/code/boost_toy.py}로 다시 계산할 수 있다. 로그 손실이 0.637 \ensuremath{\to} 0.448 \ensuremath{\to} 0.297로 매번 줄어든다.}
\Trap{나무 하나가 정답을 다 맞히게 하면(학습률 1, 깊은 나무) 학습 데이터의 우연까지 외운다. 그래서 작은 걸음(학습률 0.05)으로 천천히 쌓고, 검증 데이터에서 더 나아지지 않으면 멈춘다.}
\end{workout}

\paragraph{뉴턴 단계란.} 잎 값 = 잔차 합 \ensuremath{\div} 곡률 합. 잔차는 ``어느 방향으로 얼마나 틀렸나'', 곡률 $p(1-p)$는 ``그 점수를 바꾸면 확률이 얼마나 민감하게 변하나''다. $p$가 0.5 근처면 조금만 움직여도 되고, 0이나 1 근처면 크게 움직여야 한다.

\section{실제 설정과 학습 과정}
\begin{figure}[H]\centering
\begin{tikzpicture}[node distance=5mm,box/.style={draw=black!70,rounded corners=1.2mm,align=center,font=\small,minimum height=11mm,text width=27mm}]
\node[box] (a) {Y년 1월 1일까지\\결과 확정 대출};
\node[box,right=of a] (b) {특성 101개\\+ 정답(부도)};
\node[box,right=of b] (c) {LightGBM\\나무 최대 2000개};
\node[box,right=of c] (d) {시드 5개\\$p$ 평균};
\draw[-{Stealth},black!60] (a)--(b); \draw[-{Stealth},black!60] (b)--(c); \draw[-{Stealth},black!60] (c)--(d);
\end{tikzpicture}
\caption{연도 Y의 대출에 $p$를 붙이는 과정(시점 학습).}
\end{figure}

\begin{itemize}[leftmargin=*]
\item \textbf{학습 데이터:} 평가할 대출이 2016년 발행이면, 2016년 1월 1일까지 결과가 확정된 대출만 쓴다(부도는 마지막 납부 + 5개월에 확정, 3장). 2013년 모델은 11,345건, 2016년 모델은 156,153건으로 배운다.
\item \textbf{설정(고정, 결과를 보고 고르지 않음):} 학습률 0.05, 나무당 잎 63개, 행 80\%·\allowbreak{}열 60\% 표본추출, 잎당 최소 200건, L2 규제 5.
\item \textbf{조기 종료:} 학습 데이터의 15\%를 떼어 두고, 그 15\%의 로그 손실이 50번 연속 나아지지 않으면 나무 쌓기를 멈춘다.
\item \textbf{시드 평균:} 표본추출과 조기 종료용 분할이 시드에 따라 달라지므로 시드 5개로 모델 5개를 만들어 $p$를 평균한다.
\end{itemize}

\section{무엇이 $p$를 움직이나}
\begin{table}[H]\centering\small
\begin{tabular}{@{}rlll@{}}\toprule
순위 & 특성 & 방향 & 비중(SHAP)\\\midrule
1 & 금리 & 높을수록 위험 & 15.1\%\\
2 & 연소득 & 높을수록 안전 & 8.0\%\\
3 & 부채비율(DTI) & 높을수록 위험 & 6.1\%\\
4 & 거주 주 & 범주 & 4.1\%\\
5 & 가장 오래된 신용 계좌 시점 & 오래될수록 안전 & 3.1\%\\
6 & 만기 & 비선형 & 2.8\%\\
7 & 대출 목적 & 범주 & 2.5\%\\
8 & 대출액 & 클수록 위험 & 2.5\%\\\bottomrule
\end{tabular}
\caption{LightGBM이 $p$를 계산할 때 많이 쓴 특성(17단계).}
\end{table}
금리가 1위인 것은 자연스럽다. 금리는 LC가 이미 위험을 판단해 붙인 가격이기 때문이다. 중요한 것은 금리 다음에 소득, 부채비율, 신용 이력 같은 ``갚을 능력'' 정보가 따라온다는 점이다. 이것이 7장의 ``등급 안 정보''로 이어진다.

\section{왜 믿을 수 있나: 여섯 가지 근거}
\subsection*{근거 1. 순서가 맞는다(신뢰도 표)}
2013--16년 36개월 대출 52만 건에 시점 학습 $p$를 붙이고, $p$ 크기 순으로 10등분했다.

\begin{table}[H]\centering\small
\begin{tabular}{@{}rrrr@{}}\toprule
분위 & 평균 $p$ & 실제 부도율 & 평균 금리\\\midrule
1(가장 안전) & 2.6\% & 3.8\% & 7.3\%\\
2 & 4.2\% & 6.9\% & 9.2\%\\
3 & 5.5\% & 9.1\% & 10.2\%\\
4 & 6.9\% & 11.0\% & 11.0\%\\
5 & 8.3\% & 13.2\% & 11.6\%\\
6 & 10.1\% & 14.6\% & 12.2\%\\
7 & 12.5\% & 16.5\% & 12.8\%\\
8 & 15.7\% & 19.0\% & 13.6\%\\
9 & 20.8\% & 23.0\% & 14.8\%\\
10(가장 위험) & 33.2\% & 30.6\% & 17.2\%\\\bottomrule
\end{tabular}
\caption{$p$ 10분위별 실제 부도율(36단계). 분위마다 약 5.2만 건.}
\end{table}
실제 부도율이 한 칸도 거꾸로 가지 않고 3.8\%에서 30.6\%까지 오른다. 가장 위험한 10\%는 가장 안전한 10\%보다 8배 자주 부도났다.

\subsection*{근거 2. 같은 등급 안에서도 가른다}
등급이 이미 위험을 반영하므로, $p$가 등급만 따라 한다면 쓸모가 적다. 등급마다 $p$를 셋으로 나눴다.
\begin{table}[H]\centering\small
\begin{tabular}{@{}lrrrr@{}}\toprule
등급 & $p$ 낮음 부도율 & $p$ 중간 & $p$ 높음 & 등급 안 AUC\\\midrule
A & 3.1\% & 5.1\% & 7.9\% & 0.62\\
B & 8.4\% & 11.6\% & 15.2\% & 0.58\\
C & 14.5\% & 18.3\% & 23.4\% & 0.58\\
D & 19.8\% & 24.9\% & 30.6\% & 0.57\\
E & 23.7\% & 31.8\% & 39.6\% & 0.59\\\bottomrule
\end{tabular}
\caption{같은 등급 안에서 $p$로 나눈 실제 부도율(36단계).}
\end{table}
같은 C등급인데 $p$가 낮은 쪽은 14.5\%, 높은 쪽은 23.4\%가 부도났다. 등급이 놓친 차이를 $p$가 잡는다.

\subsection*{근거 3. 미래를 보지 않았다}
연도별 AUC는 2013년 0.631, 2014년 0.663, 2015년 0.671, 2016년 0.684다(0.5는 동전 던지기). 같은 대출로 무작위 분할을 하면 약 0.70(2007--20년 전체 종결 대출로 하면 0.73)으로 이보다 높지만, 그때 알 수 있던 정보만 쓴 \textbf{실제로 가능한 점수}다(3장). 학습 데이터가 늘수록 높아지는 것도 자연스럽다.

\subsection*{근거 4. 우연이 아니다}
\begin{itemize}[leftmargin=*]
\item \textbf{정답 섞기(위약):} 학습 라벨을 무작위로 섞어 만든 가짜 모델은 선별 가치가 0 근처였다. 평가 코드가 가짜 성과를 만들지 않는다.
\item \textbf{시드:} 시드를 바꿔도 결론이 같았다. 시드 차이(최대 0.09\%p)는 효과(약 1\%p)보다 훨씬 작다.
\item \textbf{발행월 부트스트랩:} 선별 가치 1.03\%p의 95\% 구간이 0.95--1.11로 0에서 멀다.
\end{itemize}

\subsection*{근거 5. 특정 모델의 산물이 아니다}
XGBoost는 같은 선별 가치(1.04\%p)를 냈고, 랜덤포레스트·\allowbreak{}로지스틱도 방향이 같았다(23--25단계). 결측 처리 9종, 피처 6종을 바꿔도 거의 같았다. 금리·\allowbreak{}등급을 빼고 학습한 모델도 AUC가 0.01 정도만 낮았다(34단계 F).

\subsection*{근거 6. 돈으로 이어진다}
$p$로 만든 기대수익 점수로 고르면, 실제로 들어온 돈으로 잰 수익률이 무작위 선택보다 높았다(6, 7장). 확률을 잘 맞히는 것과 돈을 버는 것은 다른 문제인데, 둘 다 확인했다.

\section{믿으면 안 되는 부분}
\begin{warnbox}[$p$의 크기는 틀린다]
\begin{itemize}[leftmargin=*]
\item 2013년 모델은 평균 $p$ 6.4\%를 냈지만 실제 부도율은 12.2\%였다.
\item \textbf{주원인은 학습 표본의 구성이다}(외부 검증의 분해). 학습 데이터에 들어가는 직전 해 발행분은 1년 안에 끝난 대출뿐인데, 그중 대부분이 조기상환이라 부도율이 최종의 절반 수준이다. 모델은 ``요즘 대출은 안전하다''고 배운다. 신용 환경 악화(추세)는 부차적이다.
\item 3장의 종결 편향(발행 후 1--3년 된 대출만 보면 부도가 과대 대표)과 방향이 반대로 보이지만 둘 다 사실이다. 나이가 1년 미만인 끝난 대출은 조기상환이, 1--3년인 끝난 대출은 부도가 과대 대표된다. $p$를 낮추는 것은 앞의 것이다.
\item 신뢰도 표의 아래쪽 분위에서 평균 $p$가 실제보다 1--4\%p 낮다.
\item 확률 보정(보정 3종)을 시도했지만 고치지 못했다(9단계). 과거로 만든 보정은 추세 변화를 따라가지 못한다.
\end{itemize}
\end{warnbox}

\begin{whybox}[그런데 왜 결과가 괜찮았나]
\textbf{직접 근거는 두 가지다.} (1) 부도확률 순(V2)은 $p$의 순서만 쓰므로, 크기가 모두 같은 비율로 틀리면 결과가 전혀 바뀌지 않는다. (2) 기대수익 순(V1)은 $s=(1-p)r-p\cdot\mathrm{LGD}-tb$에 금리 $r$이 함께 들어가 $p$의 크기가 순서를 일부 바꾼다. 그래서 $p$를 실제 부도율 수준으로 맞춰 다시 골라 봤고, 승인 대출의 88--92\%가 같았으며 전략 순위도 그대로였다.

\textbf{부수적 관찰(구조적 설명은 아님).} 점수에서 $p$는 손실률과 곱해진다. 손실률 추정(그때까지 부도 난 대출의 평균)은 실제보다 크게 나와 $p$의 과소추정과 같은 방향으로 맞물렸다. 외부 검증에 따르면 2013년에는 기대손실이 거의 맞았고, 이후 연도에는 12--75\% 크게 봤다. 이것은 결과를 사후에 본 관찰이지, 언제나 상쇄된다는 보장이 아니다.

\textbf{언제 문제가 되나.} 사람마다 틀린 정도가 달라 순서까지 망가질 때다. CatBoost가 2016년에 $p$를 평균 3.6\%로 낸 경우(실제 15.9\%)에는 성과가 무너졌다(24단계).
\end{whybox}

\begin{career}
``모델 확률을 왜 믿나요?''라는 질문에는 AUC 한 숫자로 답하지 말고, 신뢰도 표(순서), 등급 안 분리(추가 정보), 위약 검정(우연 아님), 다른 모델 비교(모델 독립), 돈으로의 연결(실제 수익)을 차례로 보여 주면 된다. 그리고 ``크기는 틀리지만 순위로만 쓴다''고 한계를 먼저 말하면 신뢰가 오히려 올라간다.
\end{career}
